Bei gleicher Temperatur haben die Moleküle beider Gase dieselbe mittlere kinetische Energie. Wegen $\bar v = \sqrt{8kT/(\pi m)}$ sind die leichteren Propanmoleküle schneller. Mit den molaren Massen (FoTaBe S.~228)
\begin{align}
 M_\text{P} &= \MPrF = \MPrg \\
 M_\text{B} &= \MBuF = \MBug
\end{align}
wird
\begin{align}
 \frac{\bar v_\text{P}}{\bar v_\text{B}} &= \vqF = \sqrt{\frac{\MBug}{\MPrg}} = \vq
\end{align}
Die Propanmoleküle sind also im Mittel um
\begin{align}
 \frac{\Delta\bar v}{\bar v_\text{B}} &= \dvrF = \dvr \approx \result{\dvrP{0}{2}}
\end{align}
schneller.

Bemerkung: Rechnet man stattdessen mit den Dichten bei \thNO\ und Normaldruck aus der Tabelle (FoTaBe S.~187), $\rho_\text{B} = \rBu$ und $\rho_\text{P} = \rPr$, und mit $p = \frac{1}{3}\rho\overline{v^2}$, erhält man $\vqrF = \vqr$, also $\dvrrP{0}{2}$. Der Unterschied kommt daher, dass Butan bei \thNO\ knapp über seiner Siedetemperatur ist und sich nicht mehr wie ein ideales Gas verhält: Seine Dichte ist grösser, als die ideale Gasgleichung ergibt.
